\section{Goldbach rings}
\label{sec:setup}

Let $\A$ be a nonempty subset of the integers bounded below, and set
\[
 R=R_\A=\kk[x_a\mid a\in\A].
\]
We write $R$, $I$, and $G$ when the support $\A$ is understood.
Our main example is the set of all positive odd primes. As throughout
the paper, $\kk$ is a field of characteristic $\ell$, with $\ell=0$
or $\ell$ prime. The estimates below hold over every field; individual
Betti numbers can nevertheless depend on the characteristic. For example,
for the full binary odd-prime ring,
$(\beta_{5,135,7},\beta_{6,135,7})$ equals $(81,9)$ outside
characteristic two and $(82,10)$ in characteristic two
\cite[Theorem~6.3]{CompanionTopology}; this example is not used in the
proofs below. Integral forms over $\mathbb Z$ carry an
explicit subscript $\mathbb Z$.

\begin{definition}
For $k\ge2$, the \emph{$k$-ary Goldbach ring} supported on $\A$ is
the $\kk$-algebra $G=G_\A=R/I$, where
\begin{equation}\label{setup:relation-ideal}
 I=I_\A=
 \bigl\langle x_{a_1}\cdots x_{a_k}-x_{b_1}\cdots x_{b_k}
 \mid a_1+\cdots+a_k=b_1+\cdots+b_k\bigr\rangle .
\end{equation}
All the labels $a_i,b_i$ belong to $\A$.
\end{definition}
When $k=2$ we call $G$ a binary Goldbach ring, and when $k=3$ a
ternary Goldbach ring. The arity $k$ specifies how many factors a
single relation may replace.
Multiplying such a relation by a monomial keeps all the other factors
fixed. Thus, in monomial degree $d\ge k$, the quotient remembers which
representations can be connected by replacements of $k$ factors at a
time. For $d<k$ no defining relation is present, while in degree $k$
all monomials with the same sum are identified. This distinguishes a
fixed-arity Goldbach ring from the toric quotient, which identifies
equal sums in every monomial degree.

We keep $k$ for the arity, $n$ for the Goldbach degree, and $d$ for the
monomial degree throughout. Statements involving binary moves will
explicitly concern $k=2$; the definitions and grading constructions
below apply to every fixed $k\ge2$.

The labels need only be bounded below. The next observation lets us
pass to positive labels when a proof needs a positive grading, without
changing the ring or its homological invariants. It also records how
the original additive weights change under that passage.

\begin{proposition}[Affine invariance]
\label{setup:affine-invariance}
Fix integers $\nu>0$ and $N$, and put
\[
 \A'=\nu\A+N=\{\nu a+N\mid a\in\A\}.
\]
For each fixed $k\ge2$, the bijection $a\mapsto \nu a+N$ induces
isomorphisms
\[
 R_\A\cong R_{\A'},\qquad G_\A\cong G_{\A'} .
\]
Under the bigrading $\deg x_a=(a,1)$ introduced below, bidegree
$(n,d)$ becomes $(\nu n+Nd,d)$.
\end{proposition}
\begin{proof}
Write $a'=\nu a+N$ and send $x_a$ to $x_{a'}$. If
$\sum_{j=1}^k a_j=\sum_{j=1}^k b_j$, then
\[
 \sum_{j=1}^k a'_j
 =\nu\sum_{j=1}^k a_j+kN
 =\nu\sum_{j=1}^k b_j+kN
 =\sum_{j=1}^k b'_j .
\]
Thus the isomorphism $R_\A\to R_{\A'}$ sends $I_\A$ into $I_{\A'}$.
Its inverse, induced by $a'\mapsto(a'-N)/\nu$, sends $I_{\A'}$
into $I_\A$. A monomial with $d$ factors and sum of labels $n$
acquires sum $\nu n+Nd$.
\end{proof}

In particular, every support bounded below is affinely equivalent
to a support contained in the positive integers. Whenever we form
support cutoffs, we first make this translation and work with positive
labels. For every integer $t$, write
\[
 \A_t=\A\cap[1,t].
\]
Thus $\A_t$ includes its upper endpoint and is empty if $t<1$.
We use $n$ for Goldbach degree and $m$ for a truncation parameter:
$\A_n$ and $\A_m$ apply the same cutoff operation, with the index
indicating its role in the argument. In particular, they are equal
when $n=m$.

\begin{remark}[Monomial notation]
For a finite multiset $\boldsymbol a=\{a_1,\ldots,a_d\}$ in $\A$, write
$x_{\boldsymbol a}=x_{a_1}\cdots x_{a_d}$; repetitions are allowed and
the empty multiset represents $1$. We use parentheses
$\boldsymbol a=(a_1,\ldots,a_d)$ when the order matters. For a finite
ordered subset $\sigma=\{a_1<\cdots<a_d\}$, put
$x_\sigma=\prod_{a\in\sigma}x_a$.

Equivalently, let $\mathcal F$ consist of the functions
$\alpha:\A\to\mathbb Z_{\ge0}$ with finite support
$\operatorname{supp}\alpha=\{a:\alpha(a)\ne0\}$. Then
\[
 x^\alpha=\prod_{a\in\operatorname{supp}\alpha}x_a^{\alpha(a)}
\]
ranges over the monomials of $R$. For instance, $\alpha(3)=2$ and
$\alpha(11)=4$, with all other values zero, represents $x_3^2x_{11}^4$.
Every polynomial is a finite sum of these monomials with coefficients
in $\kk$, uniquely after like terms are combined. These two notations
let us pass between additive representations and algebraic expressions.
\end{remark}

\subsection{Bigraded algebra structures on \texorpdfstring{$R$}{R} and \texorpdfstring{$G$}{G}}
There is a natural bigraded $\kk$-algebra structure underlying our setup
which organizes the algebraic and homological objects we study. For
$x_{\boldsymbol a}=x_{a_1}\cdots x_{a_d}$, set
\[
 \deg_1(x_{\boldsymbol a})=\sum_{j=1}^d a_j,
 \qquad \deg_2(x_{\boldsymbol a})=d.
\]
We call the first grading the \emph{Goldbach degree}, or additive weight,
and the second the \emph{monomial degree}. Homological degree is
separate from these two internal gradings. For the extension to Koszul
terms, see \hyperref[setup:koszul-grading]{the degree convention in
Section~\ref*{sec:koszul}}.

Let $R_{n,d}$ be the $\kk$-span of the monomials of bidegree $(n,d)$,
and denote their set by $\mathcal M_{n,d}$. Then
\[
 R=\bigoplus_{n\in\mathbb Z,\ d\ge0}R_{n,d},\qquad
 R_{n,d}R_{n',d'}\subseteq R_{n+n',d+d'}.
\]
Each piece is finite dimensional. Indeed, writing $b=\min\A$, every
label in a monomial of degree $d>0$ lies between $b$ and
$n-(d-1)b$.
The dimension of $R_{n,d}$ counts the unordered expressions of $n$ as
a sum of $d$ elements of $\A$, with repetitions permitted. For positive
supports these are the corresponding unordered partitions. In monomial degree
zero, $R_{0,0}=\kk$ and all other pieces vanish.

The generators of $I$ are bihomogeneous, so the quotient
inherits a bigraded algebra structure
\[
 G=\bigoplus_{n,d}G_{n,d},\qquad
 G_{n,d}G_{n',d'}\subseteq G_{n+n',d+d'}.
\]

\begin{example}[Goldbach degree $23$]
\label{setup:weight23}
Take $\A$ to be the positive odd primes and let $G$ be the binary
Goldbach ring. The nonzero pieces of $R$ in weight $23$ are
\begin{align*}
 R_{23,1}&=\kk x_{23},\\
 R_{23,3}&=\kk x_3^2x_{17}\oplus\kk x_3x_7x_{13}
          \oplus\kk x_5^2x_{13}\oplus\kk x_5x_7x_{11},\\
 R_{23,5}&=\kk x_3^4x_{11}\oplus\kk x_3^3x_7^2
          \oplus\kk x_3^2x_5^2x_7\oplus\kk x_3x_5^4,\\
 R_{23,7}&=\kk x_3^6x_5.
\end{align*}
For instance, the monomials in $R_{23,5}$ correspond to the four
unordered partitions
\[
 (3,3,3,3,11),\quad(3,3,3,7,7),\quad
 (3,3,5,5,7),\quad(3,5,5,5,5).
\]
In $G$, the four cubic monomials represent the same class, since
\begin{equation}\label{setup:weight23-cubic-moves}
 x_3^2x_{17}=x_3x_7x_{13}=x_5^2x_{13}=x_5x_7x_{11}.
\end{equation}
Each step applies one binary Goldbach move, using the pair sums
$3+17=7+13$, $3+7=5+5$, and $5+13=7+11$ in order. Similarly,
\[
 x_3^4x_{11}=x_3^3x_7^2=x_3^2x_5^2x_7=x_3x_5^4,
\]
using $3+11=7+7$ and then $3+7=5+5$ twice. Thus these two fibers
are connected by binary moves. Their common classes are nonzero,
since evaluation at $x_p=1$ descends to $G$ and sends every monomial
to $1$. Together with the singleton fibers in monomial degrees one and seven,
this proves
\[
 \dim_\kk G_{23,d}=1\quad(d=1,3,5,7).
\]
All other monomial-degree components in weight $23$ vanish.
\end{example}

For the binary ring on the odd primes, the same method works in larger
fibers. For example, in $G_{36,4}$,
\[
 x_3^2x_{11}x_{19}
 =x_3x_5x_{11}x_{17}
 =x_5x_7x_{11}x_{13}
 =x_5^2x_{13}^2 .
\]
The three moves use $3+19=5+17$, $3+17=7+13$, and $7+11=5+13$.
To count all the classes in a given bidegree we must determine which
partitions can be joined by such paths.

\subsection{The associated toric ideal}
Let $S$ be the $\kk$-algebra
\[
 S=\kk[u^a v\mid a\in\A]\subseteq\kk[u,u^{-1},v].
\]
For a positive support we can take the ambient ring to be $\kk[u,v]$.
Giving $u$ degree $(1,0)$ and $v$ degree $(0,1)$ makes
\[
 \phi:R\twoheadrightarrow S,\qquad x_a\longmapsto u^a v,
\]
a surjective bigraded $\kk$-algebra map. Its kernel $J$ is the
associated toric ideal. For the $k$-ary ideal $I$, each defining generator maps to zero:
its two terms have image $u^n v^k$, where $n$ is their common sum.
Thus $I\subseteq J$ for every fixed arity, and $\phi$ induces a
surjection $\overline\phi:G\to S$. In each represented bidegree
$(n,d)$, the target has the single basis monomial $u^n v^d$.

The map $\overline\phi$ forgets how representations are connected:
it retains only their common sum and number of factors. Its kernel
therefore measures the difference between allowing replacements of
exactly $k$ factors and allowing equal-sum relations of every arity.

To describe that kernel, take the monomials in $\mathcal M_{n,d}$ as
vertices. Join two distinct vertices when they have the form
$x^\alpha x_{a_1}\cdots x_{a_k}$ and
$x^\alpha x_{b_1}\cdots x_{b_k}$ with
$\sum_{j=1}^k a_j=\sum_{j=1}^k b_j$. The monomial $x^\alpha$ is
held fixed. Write $\Gamma_{n,d}$ for this graph, with the arity
understood from the specified ring. In the binary case $k=2$, its
edges are precisely the pair replacements used in the weight-$23$
example.

\begin{proposition}[Component bases]
\label{setup:components}
For a fixed arity $k\ge2$, one monomial from each connected
component of $\Gamma_{n,d}$ gives a basis of $G_{n,d}$ over $\kk$.
\end{proposition}
\begin{proof}
In bidegree $(n,d)$ the ideal is spanned by differences of the
endpoints of edges, the monomial multiples of the defining
binomials. Summing along a path gives the difference of its endpoints.
The map from the free module on vertices to the free module on
components sends each monomial to its component. Its kernel is
spanned by these differences, giving the asserted basis. The same
proof works for any $k$.
\end{proof}

Consequently $I=J$ if and only if each nonempty fiber is connected,
equivalently if and only if $\dim_\kk G_{n,d}\le1$ in every
bidegree. In that case $\overline\phi$ identifies $G$ with the
integral domain $S$. For the binary ring on the full odd-prime support, the converse
also follows from Theorem~\ref{depth:prime-pair-group}. Goldbach degree and
monomial degree identify the relation group with the subgroup generated by
$(p,1)$, so $\overline\phi$ becomes an isomorphism after the variables
are inverted. If $G$ is a domain it embeds in that
localization, since each monomial has nonzero image in $S$.
For general supports the converse can fail. For the binary ring on $\A=\{3,5,11\}$
the pair sums are distinct, hence $I=0$, although
$x_3^3x_{11}-x_5^4\in J$.

\begin{example}
Take the binary ring on the odd primes $\A$ congruent to $1$ modulo $4$. Then
$x_5^2x_{29}$ and $x_{13}^3$ both map to $u^{39}v^3$, but no
binary move connects them. Indeed $13+13$ is the only pair of
labels in $\A$ with sum $26$, so $x_{13}^3$ is an isolated vertex.
The other two triples of weight $39$ are joined by
$x_5^2x_{29}=x_5x_{17}^2$ in $G$. Thus $\dim_\kk G_{39,3}=2$.
\end{example}

Every nonempty binary fiber of monomial degree two is connected, since one
move joins any two decompositions of the pair sum. For a binary ring on a finite
support $\A$, write $G_{*,d}=\bigoplus_n G_{n,d}$, where $*$ denotes
the direct sum over all Goldbach degrees. Then
\begin{equation}\label{setup:low-hilbert}
 \dim_\kk G_{*,1}=|\A|,\qquad \dim_\kk G_{*,2}=|\A+\A|.
\end{equation}
For the toric quotient, write $S_{*,d}=\bigoplus_n S_{n,d}$.
Its basis consists of the monomials $u^n v^d$ with $n\in d\A$, where
$d\A$ denotes the $d$-fold sumset. Thus, for every finite support,
\[
 \dim_\kk S_{*,d}=|d\A|.
\]
This identifies iterated sumset cardinalities with the Hilbert function
of a standard graded algebra, the correspondence used by
Eliahou--Mazumdar \cite{EliahouMazumdar22} and Elias \cite{Elias22}.

\paragraph{The curve defined by the toric ideal.}
The toric quotient provides a geometric comparison for the Goldbach
ring. Its monomials are indexed by attained bidegrees rather than by
move components, so its Hilbert function is simpler. This comparison
does not identify $G$ with $S$: connectivity, cancellation, and depth
must still be checked in the relevant ring. On the odd-prime support
the toric quotient has a particularly explicit curve description,
independent of which arity was used to define $G$.

For the odd primes write $y=u^3v$, $t=u^2$ and $e_p=(p-3)/2$.
Then $S=\kk[yt^{e_p}:p\in\Podd]$. Write $S^{<m}$ for the subalgebra
generated by the $x_p=yt^{e_p}$ with $p\in\A_m$. For $m\ge5$,
the first two generators are $y$ and $yt$, so the fraction field
of $S^{<m}$ is $\kk(y,t)$. Thus $\operatorname{Spec}S^{<m}$ is a
two-dimensional affine cone. With the monomial grading its projective
curve is the closure of $t\mapsto[1:t:\cdots:t^{e_P}]$, where $P$
is the largest prime at most $m$ and the omitted coordinates are the
other $t^{e_p}$ in increasing order. This rational monomial curve is
locally Cohen--Macaulay, since its local rings at closed points are
one-dimensional Noetherian domains.

The finite exponent sets $\{e_p:p\in\A_m\}$ place these curves in the
sumset framework of Elias \cite{Elias22}. Relations between sumset
structure, Castelnuovo--Mumford regularity, and the Cohen--Macaulay
property of their homogeneous coordinate rings are developed by
Gimenez--Gonz\'alez-S\'anchez \cite{GimenezGonzalez23} and
Hoa--Tien \cite{HoaTien26}. The recent preprint
\cite{ChaudharyKhandelwal26} further studies the difference between
sumset regularity and Castelnuovo--Mumford regularity.

The cone can have depth one. At $m=42$, $x_3,x_{41}$ are homogeneous
parameters since $x_p^{19}=x_3^{19-e_p}x_{41}^{e_p}$ in $S^{<42}$.
The class of $x_5x_7$ modulo $x_3$ is nonzero: dividing by $x_3$ would
require a coefficient of monomial degree one and Goldbach degree $9$.
Yet $x_{41}x_5x_7=x_3x_{13}x_{37}$, so this pair of parameters is
not a regular sequence. The domain $S^{<42}$ therefore has depth
one at the vertex over every field.

On the full support $S$ still has dimension two. Any finite strict
chain of prime ideals contracts strictly to a sufficiently large
finite cutoff: include one witness to each strict inclusion. For the
lower bound, the ideal generated by the $x_p$ with $p>3$ has quotient
$\kk[x_3]$ and lies strictly between $(0)$ and $S_+=\bigoplus_{d>0}S_{*,d}$.
For $q>p>3$ we have $x_p^{e_q}=x_3^{e_q-e_p}x_q^{e_p}$, giving
$\sqrt{x_3 S}=S_+$. Thus $D(x_3)$ is the complement of the vertex,
and $S[x_3^{-1}]=\kk[y^{\pm1},t]$. Taking its monomial-degree-zero part gives
$\operatorname{Proj}S=D_+(x_3)\cong\mathbb A^1_\kk$.
At the vertex $S_+/S_+^2$ has one basis element for every odd prime.
The local ring there is non-Noetherian, and the maximum length of a
regular sequence in its maximal ideal is one. Indeed, $x_3$ is regular.
For any two nonzero nonunits $f,g$ in that local ring, some powers
satisfy $f^N=x_3 a$ and $g^M=x_3 b$. Then $a\notin(f^N)$ but
$a g^M=b f^N$, ruling out a regular sequence $f,g$ since powers of
a regular sequence are regular.

Finally, affine invariance supplies a positive grading without changing
the original bidegrees. Translating every label by $1-b$ gives the
auxiliary weight
\begin{equation}\label{setup:positive-weight}
 w(n,d)=n-(b-1)d,\qquad w(x_a)=a-b+1\ge1.
\end{equation}
There are finitely many variables, and consequently finitely many
monomials, in every bounded auxiliary-weight range. This connected,
locally finite positive grading will allow us to construct minimal
resolutions even when $\A$ is infinite.
